\chapter{The Long-Run Model}\label{ch:b1:long-run-model}

Output per worker in the rich countries is several times what it was in
1950, yet the share of national income paid to labour has moved within a
narrow band: in the latest year of data (2023, and 2022 for Japan) it was
between $0.568$ and $0.637$ in the United States, Japan, the United Kingdom,
France and Germany. % figdata: labsh_points
Why should the split of income be steadier than income itself? And what
sets the interest rate, the price at which a country's saving is lent to
those who invest? This chapter builds the first macroeconomic model of the
book: an economy in the long run, with prices free to adjust, the factors
fully employed and money left aside. It is the model the short run of
Part~VI departs from.

\begin{openingchart}{Labour's share of GDP, five countries, 1950--2023. Years
that the table fills with a country's first or last observed value are left
out: Germany starts in 1980, Japan and the United Kingdom in 1970, and
Japan ends in 2022.}{Penn World Table 11.0 (labsh)}
\begin{tikzpicture}
\begin{axis}[omtime, xlabel={Year}, ylabel={Labour share of GDP}, xmin=1950, xmax=2023, ymin=0.45, ymax=0.75,
  width=0.72\linewidth, height=0.42\linewidth, legend style={font=\footnotesize, at={(0.5,-0.24)}, anchor=north,
  legend columns=5}]
  \addplot+[thick, no marks] table[x=year, y=fra, col sep=comma] {figdata/bachelor-1/19-long-run-model/out/labsh.csv};
  \addplot+[thick, no marks] table[x=year, y=deu, col sep=comma] {figdata/bachelor-1/19-long-run-model/out/labsh.csv};
  \addplot+[thick, no marks] table[x=year, y=jpn, col sep=comma] {figdata/bachelor-1/19-long-run-model/out/labsh.csv};
  \addplot+[thick, no marks] table[x=year, y=gbr, col sep=comma] {figdata/bachelor-1/19-long-run-model/out/labsh.csv};
  \addplot+[thick, no marks] table[x=year, y=usa, col sep=comma] {figdata/bachelor-1/19-long-run-model/out/labsh.csv};
  \legend{France, Germany, Japan, United Kingdom, United States}
\end{axis}
\end{tikzpicture}
\end{openingchart}

The shares are stable but not constant. Japan's rose from $0.508$ in 1970,
its first observed year, to $0.61$ in 1980; the American share was $0.628$ in 1950 and $0.568$ in 2023.
% figdata: labsh_points
Since 1980 the share has fallen in four of the five countries, by $1.3$
points a decade in Germany and $0.78$ in the United States, while it rose by
$0.75$ points a decade in the United Kingdom. % figdata: labsh_trend
A model with fixed shares explains the band; it does not explain the
drift.

\begin{recall}
The aggregate \omterm{def:b1:technology:production-function}{production function} and capital's share, saving and
investment, \omterm{def:b1:long-run-model:crowding-out}{crowding out} and the \omterm{def:b1:long-run-model:long-run}{long-run model} in words were met at
school. Euler's theorem is \cref{thm:b1:technology:euler}; profit
maximisation is in \cref{ch:b1:profit-maximisation}; \omterm{def:b1:factor-markets:factor-share}{factor shares} and
their measurement are in \cref{ch:b1:factor-markets}.
\end{recall}

\section{Output from the factors}\label{sec:b1:long-run-model:output}

The economy produces one good with capital $K$ and labour $L$:
$Y = F(K, L)$, with \omterm{def:b1:technology:returns-to-scale}{constant returns to scale}. In the long run of this
chapter, $K$ and $L$ are given: the capital stock was built in the past, and
the labour force is set by population and by households' choice to work.
Output is therefore given too, $\bar Y = F(\bar K, \bar L)$. What the model determines is
how that output is divided between the factors and between its uses.

\begin{definition}[Long-run equilibrium]\label{def:b1:long-run-model:long-run}
\index{long-run equilibrium}\index{long-run model}
A \emph{long-run \omterm{def:b1:models-and-data:equilibrium}{equilibrium}} of the economy is a state in which all
prices --- wages, the \omterm{def:b1:factor-markets:rental-price}{rental price of capital}, the interest rate --- have
adjusted so that the factors are fully employed and output equals the
demand for it.
\end{definition}

\section{Factor prices}\label{sec:b1:long-run-model:factors}

Competitive firms hire each factor until its \omterm{def:b1:technology:marginal-product}{marginal product} equals its
real price (\cref{ch:b1:profit-maximisation}): $w = F_L(K, L)$ and
$R/P = F_K(K, L)$. With constant returns, by Euler's theorem,
\[
F_K K + F_L L = Y:
\]
paying each factor its \omterm{def:b1:technology:marginal-product}{marginal product} exhausts output, and there is no
pure profit left. For the Cobb--Douglas $Y = AK^{\alpha}L^{1-\alpha}$, $F_L L/Y = 1 - \alpha$ and
$F_K K/Y = \alpha$: the shares are constant, whatever $K$, $L$ and $A$.

\begin{omfigure}
\begin{tikzpicture}
\begin{axis}[omrecord, grid=none, xlabel={Labour $L$ (million workers)}, ylabel={Real wage},
  xmin=0, xmax=8, ymin=0, ymax=150, width=0.5\linewidth, height=0.4\linewidth, samples=60]
  \addplot[demand, domain=1:8] {0.65*13.6087*1000^0.35*x^(-0.35)};
  \addplot[supply] coordinates {(4.4,0) (4.4,150)};
  \addplot[guide] coordinates {(0,59.09) (4.4,59.09)};
  \node[font=\footnotesize, anchor=west] at (axis cs:6,30) {$F_L(\bar K, L)$};
  \node[font=\footnotesize, anchor=south west] at (axis cs:4.5,120) {$\bar L$};
  \node[font=\footnotesize, anchor=south] at (axis cs:2,60) {$w = 59.09$};
\end{axis}
\end{tikzpicture}

{\small The labour market of the weekend problem: the \omterm{def:b1:technology:marginal-product}{marginal product} of
labour with $\bar K = 1000$ (blue) and a fixed labour supply of 4.4 million (red).
The real wage, 59.09 thousand dalens a year, is the \omterm{def:b1:technology:marginal-product}{marginal product} at full
employment.}
\end{omfigure}

\begin{record}[The fall of the labour share]
Labour's share was long taken as a constant of macroeconomics. Measured
with care across many countries and industries, it has declined
significantly since the early 1980s; one study attributes roughly half of
the decline to the falling price of investment goods, which led firms to
substitute capital for labour (Karabarbounis and Neiman 2014).
% ledger: b1:karabarbounis2014
With Cobb--Douglas technology that cannot happen: the shares are fixed. With
an \omterm{def:b1:technology:sigma}{elasticity of substitution} above one, cheaper capital raises capital's
share; the estimates of \cref{ch:b1:technology}, mostly below one, point the
other way, and the question is open. Part of the measured fall is also a
matter of measurement: measured shares depend on how the income of the
self-employed is split between labour and capital (Gollin 2002, for
differences between countries; \cref{ch:b1:factor-markets}). % ledger: b1:gollin2002
\end{record}

\section{Uses of output}\label{sec:b1:long-run-model:uses}

In a closed economy output is consumed, invested or bought by the
government: $Y = C + I + G$. In the \omterm{def:b1:long-run-model:long-run}{long-run model}:
\begin{itemize}
  \item consumption depends on \omterm{def:b1:national-accounts-in-depth:gni}{disposable income}, $C = C(Y - T)$, with a
  marginal propensity to consume between 0 and 1;
  \item investment depends on the real interest rate, $I = I(r)$, decreasing:
  a project is worth undertaking when its return exceeds the cost of
  borrowing;
  \item government spending $G$ and taxes $T$ are set by policy.
\end{itemize}

\begin{definition}[National, private and public saving]\label{def:b1:long-run-model:saving}
\index{national saving}\index{private saving}\index{public saving}
\emph{Private saving} is \omterm{def:b1:national-accounts-in-depth:gni}{disposable income} not consumed, $Y - T - C$;
\emph{public saving} is tax revenue less government spending, $T - G$;
\emph{national saving} is their sum, $S = Y - C - G$.
\end{definition}

\section{The market for loanable funds}\label{sec:b1:long-run-model:loanable}

\begin{definition}[Loanable funds market]\label{def:b1:long-run-model:loanable-funds}
\index{loanable funds market}
The \emph{loanable funds market} is the market in which \omterm{def:b1:long-run-model:saving}{national saving} is
supplied to those who invest; its price is the real interest rate, and in
\omterm{def:b1:models-and-data:equilibrium}{equilibrium} $S = I(r)$.
\end{definition}

The identity $Y = C + I + G$ is the same statement: $Y - C - G = I$. In the long
run $Y$ is given and $C$ depends only on $Y - T$, so \omterm{def:b1:long-run-model:saving}{national saving} does not
depend on $r$, and the \omterm{def:b1:models-and-data:equilibrium}{equilibrium} interest rate is the one at which
investment demand absorbs it.

\begin{definition}[Crowding out]\label{def:b1:long-run-model:crowding-out}
\index{crowding out}
\emph{Crowding out} is the fall in private investment (or other private
spending) that follows a rise in government spending or borrowing, through
a higher interest rate.
\end{definition}

\begin{proposition}[Full crowding out]\label{prop:b1:long-run-model:crowding}
In the \omterm{def:b1:long-run-model:long-run}{long-run model}, a rise $\Delta G$ in government purchases with taxes
unchanged lowers investment by exactly $\Delta G$ and raises the real interest
rate by $\Delta G/|I'(r)|$. A cut $\Delta T$ in taxes lowers investment by $\mathrm{MPC}\cdot\Delta T$.
\end{proposition}

\begin{proof}
$S = \bar Y - C(\bar Y - T) - G$ falls by $\Delta G$ when $G$ rises; in \omterm{def:b1:models-and-data:equilibrium}{equilibrium} $I = S$, so $I$
falls by $\Delta G$, which requires $I'(r)\,\Delta r = -\Delta G$. A tax cut raises $C$ by $\mathrm{MPC}\cdot\Delta T$
and lowers $S$, hence $I$, by the same amount.
\end{proof}

\begin{omfigure}
\begin{tikzpicture}
\begin{axis}[omrecord, grid=none, xlabel={Saving and investment (billions)}, ylabel={Real interest rate (\%)},
  xmin=50, xmax=130, ymin=0, ymax=8, width=0.52\linewidth, height=0.42\linewidth]
  \addplot[demand, domain=52:132] {(132 - x)/10};
  \addplot[supply] coordinates {(92,0) (92,8)};
  \addplot[shifted] coordinates {(84,0) (84,8)};
  \addplot[guide] coordinates {(50,4) (92,4)};
  \addplot[guide] coordinates {(50,4.8) (84,4.8)};
  \node[font=\footnotesize, anchor=west] at (axis cs:110,1.7) {$I(r)$};
  \node[font=\footnotesize, anchor=south west] at (axis cs:92,7) {$S$};
  \node[font=\footnotesize, anchor=south east] at (axis cs:84,7) {$S'$};
\end{axis}
\end{tikzpicture}

{\small Loanable funds in the weekend problem: \omterm{def:b1:long-run-model:saving}{national saving} of 92 billion
(vertical, red) and investment $I = 132 - 10r$ (blue) give $r = 4\%$. Defence
spending of 8 billion more shifts saving to 84 (dashed): $r$ rises to $4.8\%$ and
investment falls by 8.}
\end{omfigure}

The model predicts a link between saving, investment and the interest rate
that is hard to see in the data, because the curves move together.

\begin{omchart}{United States: investment as a share of GDP against the real
long-term interest rate, 1950--2020; the real rate is the long-term rate
minus average CPI inflation over the previous three years.}{Jordà--Schularick--Taylor
Macrohistory Database, release 6}
\begin{tikzpicture}
\begin{axis}[omscatter, xlabel={Real long-term rate (\%)}, ylabel={Investment (\% of GDP)}, xmin=-4, xmax=10,
  ymin=16, ymax=24, width=0.6\linewidth, height=0.4\linewidth]
  \addplot[only marks, mark=*, mark size=1.3pt, omDef] table[x=real, y=inv, col sep=comma] {figdata/bachelor-1/19-long-run-model/out/us_invest.csv};
\end{axis}
\end{tikzpicture}
\end{omchart}

The cloud has no slope: the correlation is $0.11$ over 71 years, and
investment stayed between $17.98\%$ and $23.03\%$ of GDP. % figdata: us_invest_stats
That does not refute $I'(r) < 0$: when investment demand shifts out, it raises
both $I$ and $r$; when saving shifts, it moves them in opposite directions.
The points mix both, as the market for fish did in \cref{ch:b1:models-and-data}.

Evidence on \omterm{def:b1:long-run-model:crowding-out}{crowding out} therefore comes from designs that separate the
curves. One holds the business cycle and monetary policy aside by using
the interest rates expected five to ten years ahead and the budget office's
projections of future deficits: in the United States, a projected deficit
one point of GDP larger goes with long-term rates about 25 basis points
higher, and a debt one point of GDP larger with rates 3 to 4 basis points
higher (Laubach 2009). % ledger: b1:laubach2009
Interest rates rise, as the model predicts, but by an amount that does not
by itself show how much investment falls. Another design compares firms in
the same country: in China from 2006 to 2013, private firms invested less in
cities with more local public debt, which tightened their access to credit,
while the investment of state-owned firms did not change (Huang, Pagano and
Panizza 2020). % ledger: b1:hpp2020
Neither study measures \omterm{def:b1:long-run-model:crowding-out}{crowding out} as a whole: the first sees the price of
loanable funds, the second one channel and one country.

\section{Money apart}\label{sec:b1:long-run-model:money}

Nothing so far involved money. Output, the real wage, the real interest
rate and the split of output are all set by technology, the factor
supplies, preferences and fiscal policy.

\begin{definition}[Neutrality of money]\label{def:b1:long-run-model:neutrality}
\index{neutrality of money}\index{monetary neutrality}
Money is \emph{neutral} when changes in the quantity of money change
nominal variables --- prices, nominal wages --- but not real ones: output,
employment, real wages and real interest rates.
\end{definition}

In the \omterm{def:b1:long-run-model:long-run}{long-run model} money only sets the price level, the subject of
\cref{ch:b1:money-and-inflation}. In the short run wages and prices adjust
slowly, money is not \omterm{def:b1:long-run-model:neutrality}{neutral}, and output can depart from $F(\bar K, \bar L)$: that
is the subject of Part~VI.

\begin{model}{The long-run model}
\assumes Flexible prices and wages; full employment of given capital and
labour; a closed economy; competitive factor markets; consumption that
depends on \omterm{def:b1:national-accounts-in-depth:gni}{disposable income} and investment on the real interest rate.
\predicts \omterm{def:b1:factor-markets:factor-share}{Factor shares} set by technology (constant with Cobb--Douglas); a
real interest rate that clears the \omterm{def:b1:long-run-model:loanable-funds}{loanable funds market}; full \omterm{def:b1:long-run-model:crowding-out}{crowding out}
of investment by government purchases; \omterm{def:b1:long-run-model:neutrality}{neutral} money.
\tested Consistent with the shares: where they are observed, \omterm{def:b1:factor-markets:factor-share}{labour shares}
stayed within about a tenth of GDP for five decades or more (opening
chart), but they have declined since the early 1980s (Karabarbounis and
Neiman 2014). % ledger: b1:karabarbounis2014
\omterm{def:b1:long-run-model:crowding-out}{Crowding out} is tested only in part: larger projected US deficits raise
long-term interest rates (Laubach 2009), % ledger: b1:laubach2009
and local public debt cut private investment in Chinese cities (Huang,
Pagano and Panizza 2020); % ledger: b1:hpp2020
neither measures whether it is full, and the US cloud cannot separate the
curves. Long-run \omterm{def:b1:long-run-model:neutrality}{neutrality of money} is tested in \cref{ch:b1:money-and-inflation}.
\breaks In the short run (Part~VI); in open economies, where saving can
come from abroad (\cref{ch:b1:open-economy}); and the decline of \omterm{def:b1:factor-markets:factor-share}{labour
shares} since 1980 contradicts constant Cobb--Douglas shares.
\end{model}

\section*{Exercises}\label{sec:b1:long-run-model:exercises}

\begin{exercise}[$\star$]\label{exo:b1:long-run-model:prices}
With $Y = 10K^{0.3}L^{0.7}$, $K = 100$ and $L = 400$, compute output, the real wage and
the \omterm{def:b1:factor-markets:rental-price}{rental price of capital}.
\end{exercise}

\begin{exercise}[$\star$]\label{exo:b1:long-run-model:read}
From the opening chart's table, the \omterm{def:b1:factor-markets:factor-share}{labour share} was $0.508$ in Japan in
1970, its first observed year, and $0.628$ in the United States in 1950; in
the latest year it was $0.587$ in Japan (2022) and $0.568$ in the United
States (2023).
% figdata: labsh_points
Describe the change in each.
\end{exercise}

\begin{exercise}[$\star$]\label{exo:b1:long-run-model:euler-ex}
Check Euler's theorem on the numbers of
\cref{exo:b1:long-run-model:prices}.
\end{exercise}

\begin{exercise}[$\star$]\label{exo:b1:long-run-model:predict}
The government raises its purchases and leaves taxes unchanged. Predict
the effect on the real interest rate, investment and consumption in the
\omterm{def:b1:long-run-model:long-run}{long-run model}.
\end{exercise}

\begin{exercise}[$\star\star$]\label{exo:b1:long-run-model:lf}
$Y = 2000$, $T = G = 300$, $C = 100 + 0.6(Y - T)$ and $I = 700 - 20r$. Find private,
public and \omterm{def:b1:long-run-model:saving}{national saving}, the interest rate and investment.
\end{exercise}

\begin{exercise}[$\star\star$]\label{exo:b1:long-run-model:taxcut}
In \cref{exo:b1:long-run-model:lf}, taxes are cut by 50. Find the new $r$, $I$
and $C$.
\end{exercise}

\begin{exercise}[$\star\star$]\label{exo:b1:long-run-model:immig}
With the technology of \cref{exo:b1:long-run-model:prices}, immigration
raises $L$ by $10\%$. Predict the change in the real wage, the \omterm{def:b1:factor-markets:rental-price}{rental price of
capital} and the \omterm{def:b1:factor-markets:factor-share}{factor shares}.
\end{exercise}

\begin{exercise}[$\star\star$]\label{exo:b1:long-run-model:trend}
From the opening chart's table, the least-squares trend of the \omterm{def:b1:factor-markets:factor-share}{labour share}
since 1980 is $-1.3$ points a decade in Germany, $-1.15$ in France, $-1.01$ in
Japan, $-0.78$ in the United States and $+0.75$ in the United Kingdom.
% figdata: labsh_trend
By how much would each have moved over the 43 years to 2023 at that rate?
\end{exercise}

\begin{exercise}[$\star\star$]\label{exo:b1:long-run-model:brenland}
In Brenland's accounts for 2025, $Y = 393{,}264$, $C = 225{,}014$, $G = 78{,}653$ and
$I = 88{,}484$ (millions of dalens). Compute \omterm{def:b1:long-run-model:saving}{national saving} as in a closed
economy and explain the gap with investment.
\end{exercise}

\begin{exercise}[$\star\star\star$]\label{exo:b1:long-run-model:ces}
With a CES technology and $\sigma < 1$, show that a rise in $K/L$ raises labour's
share. What happens with $\sigma > 1$?
\end{exercise}

\begin{exercise}[$\star\star\star$]\label{exo:b1:long-run-model:partial}
Suppose \omterm{def:b1:long-run-model:saving}{national saving} rises with the interest rate, $S = \bar S + sr$. Derive the
effect of $\Delta G$ on $r$ and $I$, and show that \omterm{def:b1:long-run-model:crowding-out}{crowding out} of investment is
partial.
\end{exercise}

\begin{exercise}[$\star\star\star$]\label{exo:b1:long-run-model:neutral-ex}
What observation would refute the \omterm{def:b1:long-run-model:neutrality}{neutrality of money} in the long run?
\end{exercise}

\section{Problem: Brenland in the Long Run}\label{sec:b1:long-run-model:problem}

\begin{problem}[{Weekend problem --- a defence build-up}]\label{pb:b1:long-run-model:brenland}
Brenland's economy (invented) is modelled with shares close to those of its
2025 accounts. Output is $Y = AK^{0.35}L^{0.65}$ with $K = 1000$ billion dalens and
$L = 4.4$ million workers, and $Y = 400$ billion. Consumption is
$C = 36 + 0.6(Y - T)$, taxes $T = 80$, government purchases $G = 80$, and investment
$I = 132 - 10r$, with $r$ in per cent (all in billions of dalens a year).

\medskip\noindent\textbf{Part I --- Factor prices.}
\begin{enumerate}
  \item Find $A$.
  \item Compute the real wage (thousands of dalens per worker a year).
  \item Compute the \omterm{def:b1:factor-markets:rental-price}{rental price of capital} (per dalen of capital a year).
  \item Check Euler's theorem.
  \item What are labour's and capital's shares?
  \item By what factors would the wage and the rental price change if $L$ rose by $10\%$?
\end{enumerate}

\medskip\noindent\textbf{Part II --- Loanable funds.}
\begin{enumerate}[resume]
  \item Compute consumption.
  \item Compute private and \omterm{def:b1:long-run-model:saving}{public saving}.
  \item Compute \omterm{def:b1:long-run-model:saving}{national saving}.
  \item Find the real interest rate and investment.
  \item Compare the shares of $C$, $I$ and $G$ in output with Brenland's 2025 accounts: $0.572$, $0.225$ and $0.2$.
  \item What would happen to $r$ if investment demand rose by 10 at every rate?
\end{enumerate}

\medskip\noindent\textbf{Part III --- A defence build-up.} Defence purchases rise
by $2\%$ of output, financed by borrowing.
\begin{enumerate}[resume]
  \item By how much do public and \omterm{def:b1:long-run-model:saving}{national saving} change?
  \item Find the new interest rate.
  \item By how much does investment fall?
  \item If the build-up were financed by taxes instead, find $C$, \omterm{def:b1:long-run-model:saving}{national saving}, $r$ and $I$.
  \item Why does tax finance crowd out less investment?
  \item Over the years, what does lower investment do to the capital stock and to wages?
\end{enumerate}

\medskip\noindent\textbf{Part IV --- A technology gain.} $A$ rises by $10\%$, with $K$, $L$,
$T$, $G$ and investment demand unchanged.
\begin{enumerate}[resume]
  \item Find output.
  \item Find the wage and the rental price.
  \item What happens to the \omterm{def:b1:factor-markets:factor-share}{factor shares}?
  \item Find \omterm{def:b1:long-run-model:saving}{national saving} and the interest rate.
  \item Why might investment demand not stay unchanged, and which way would that move $r$?
  \item State the result: the effect of the defence build-up on the real interest rate and on investment.
\end{enumerate}
\end{problem}
